\begin{figure}[htbp]
\centering
\begin{tikzpicture}[
  vertex/.style={circle,draw=black!65,fill=white,inner sep=2pt},
  edge/.style={-{Stealth[length=1.8mm]},draw=black!65},
  lab/.style={font=\small,fill=white,inner sep=1pt}]
 \node[vertex] (j) at (0,0) {$j$};
 \node[vertex] (g) at (1.65,1) {$g$};
 \node[vertex] (k) at (1.65,-1) {$k$};
 \node[vertex] (i) at (3.3,0) {$i$};
 \draw[edge] (j) -- node[lab,above left] {$C$} (g);
 \draw[edge] (j) -- node[lab,below left] {$Q^{\mathsf T}$} (k);
 \draw[edge] (g) -- node[lab,above right] {$B$} (i);
 \draw[edge] (k) -- node[lab,below right] {$E$} (i);
 \draw[edge,line width=1pt] (g) -- node[lab,right] {$T^{\mathsf T}$} (k);
 \node at (4.3,0) {$\longrightarrow$};
 \node[vertex] (jj) at (5.3,0) {$j$};
 \node[vertex] (v) at (7.6,0) {$a$};
 \node[vertex] (ii) at (9.9,0) {$i$};
 \draw[edge] (jj) -- node[lab,above] {$C\circ Q^{\mathsf T}$} (v);
 \draw[edge] (v) -- node[lab,above] {$B\circ E$} (ii);
 \node[font=\small,below=5mm of v] {$T_{aa}$};
 \node[font=\small] at (1.65,-1.65) {two internal indices};
 \node[font=\small] at (7.6,-1.65) {one internal index};
\end{tikzpicture}
\caption{Pinching the bridge in a crossed contraction.
The left network represents
$\sum_{g,k}B_{ig}E_{ik}T_{gk}C_{gj}Q_{jk}$.
Removing the off-diagonal entries of $T$ costs at most a fixed
operator-norm constant times $\max_{g\ne k}|T_{gk}|$.
Its diagonal identifies $g=k=a$, leaving
$(B\circ E)D_{\operatorname{diag}T}(C\circ Q^{\mathsf T})$.
The diagonal label is retained.}
\label{fig:network-pinching}
\end{figure}
